\documentclass[troisieme]{fiche}
\metadonnees{3}{Le pré}{Bloc A — Raconter}

\begin{document}
\entete

\objectifs{%
\textbf{Langue} : prétérit, formes négative et interrogative ; verbes irréguliers, lot 2.\par
\textbf{Texte} : Anna Sewell, \emph{Black Beauty} (1877), ch. \textsc{i}.\par
\textbf{Durée} : 1 h — exercices 20 min, texte et questions 25 min, version 15 min.}

%% ===================================================================
\exercice[12 min]{Négation et question au prétérit}

\rappel{\textbf{Rappel.} Au prétérit, la négation et la question passent par \emph{did}, et
le verbe revient à sa base. Seul \emph{be} se passe d'auxiliaire.}

\consigne{Réécrivez la phrase comme demandé.}

\begin{anglais}
\begin{enumerate}[label=\arabic*.,itemsep=3mm,leftmargin=*]
  \item The colts learned good manners. \emph{(forme négative)}
    \lignes{1}
    \corr{The colts did not learn good manners.}
  \item He ran with the other colts. \emph{(forme interrogative)}
    \lignes{1}
    \corr{Did he run with the other colts?}
  \item His mother bit and kicked. \emph{(forme négative)}
    \lignes{1}
    \corr{His mother did not bite or kick.}
  \item The master was kind to them. \emph{(forme interrogative)}
    \lignes{1}
    \corr{Was the master kind to them?}
  \item They had great fun in the field. \emph{(forme négative)}
    \lignes{1}
    \corr{They did not have great fun in the field.}
  \item Duchess won the cup at Newmarket. \emph{(forme interrogative)}
    \lignes{1}
    \corr{Did Duchess win the cup at Newmarket?}
  \item He forgot his mother's advice. \emph{(forme négative)}
    \lignes{1}
    \corr{He did not forget his mother's advice.}
  \item The colts were nearly as large as grown-up horses. \emph{(forme négative)}
    \lignes{1}
    \corr{The colts were not nearly as large as grown-up horses.}
\end{enumerate}
\end{anglais}

\note[La faute à ne pas faire]{Une seule marque du passé par groupe verbal : \emph{did not
learn}, jamais \emph{did not learned}. De même à l'interrogatif : \emph{did he run?} et non
\emph{did he ran?}}

\note[Deux détails]{Phrase 3 : après une négation, \emph{and} devient \emph{or} (\emph{did
not bite or kick}). Phrases 4 et 8 : avec \emph{be}, pas de \emph{did} — on inverse
(\emph{was the master\ldots?}) ou on ajoute \emph{not} (\emph{were not}).}

%% ===================================================================
\exercice[8 min]{Verbes irréguliers, lot 2}
\consigne{a) Complétez.}

\begin{center}\small
\begin{tabular}{@{}lll@{}}
\textsf{base} & \textsf{prétérit} & \textsf{participe passé} \\[1.5mm]
\emph{give} & \trou{gave} & \trou{given} \\[1mm]
\emph{find} & \trou{found} & \trou{found} \\[1mm]
\emph{think} & \trou{thought} & \trou{thought} \\[1mm]
\emph{know} & \trou{knew} & \trou{known} \\[1mm]
\emph{run} & \trou{ran} & \trou{run} \\[1mm]
\emph{eat} & \trou{ate} & \trou{eaten} \\[1mm]
\emph{drink} & \trou{drank} & \trou{drunk} \\[1mm]
\emph{begin} & \trou{began} & \trou{begun} \\[1mm]
\emph{break} & \trou{broke} & \trou{broken} \\[1mm]
\emph{bring} & \trou{brought} & \trou{brought} \\
\end{tabular}
\end{center}

\consigne{b) Complétez avec un verbe du tableau, au prétérit.}
\begin{anglais}
\begin{enumerate}[label=\arabic*.,itemsep=2.5mm,leftmargin=*]
  \item While he was young he \trou{drank} his mother's milk.
  \item He \trou{ran} with the other colts all round the field.
  \item His mother \trou{gave} him her advice one day, and he never forgot it.
  \item He \trou{knew} she was a wise old horse.
\end{enumerate}
\end{anglais}

\note[Trois familles]{\emph{Think} et \emph{bring} font \emph{-ought} : même son,
orthographes proches. \emph{Drink}, \emph{begin} et \emph{run} suivent le schéma
\emph{i\,--\,a\,--\,u}, comme \emph{sing}, \emph{swim}, \emph{ring}. \emph{Give},
\emph{eat}, \emph{break}, \emph{know} ont un participe en \emph{-en}.}

%% ===================================================================
\newpage
\section{Texte}

\noindent\small\textit{Le narrateur de ce roman est un cheval. Il raconte sa vie depuis sa
naissance.}\normalsize

\begin{texte}
While I was young I lived upon my mother's milk, as I could not eat grass. In the daytime I
ran by her side, and at night I lay down close by her. When it was hot we used to stand by the
pond in the shade of the trees, and when it was cold we had a nice warm
\gl[a shed]{shed}{un abri, un appentis} near the \gl[a grove]{grove}{un bosquet}.

As soon as I was old enough to eat grass my mother used to go out to work in the daytime, and
come back in the evening.

There were six young \gl[a colt]{colts}{un poulain} in the meadow besides me; they were older
than I was; some were nearly as large as grown-up horses. I used to run with them, and had
great fun; we used to \gl[to gallop]{gallop}{galoper} all together round and round the field
as hard as we could go. Sometimes we had rather rough play, for they would frequently bite and
kick as well as gallop.

One day, when there was a good deal of kicking, my mother \gl[to whinny]{whinnied}{hennir} to
me to come to her, and then she said:

``I wish you to pay attention to what I am going to say to you. The colts who live here are
very good colts, but they are \gl[a cart-horse]{cart-horse}{un cheval de trait} colts, and of
course they have not learned manners. You have been well-\gl[to breed, bred]{bred}{élever}
and well-born; your father has a great name in these parts, and your grandfather won the cup
two years at the Newmarket races; your grandmother had the sweetest
\gl[a temper]{temper}{un caractère} of any horse I ever knew, and I think you have never seen
me kick or bite. I hope you will grow up gentle and good, and never learn bad ways; do your
work with a good will, lift your feet up well when you \gl[to trot]{trot}{trotter}, and never
bite or kick even in play.''

I have never forgotten my mother's advice; I knew she was a wise old horse, and our master
thought a great deal of her. Her name was Duchess, but he often called her Pet.
\end{texte}
\source{Anna Sewell, \emph{Black Beauty}, Londres, Jarrold, 1877, ch.~\textsc{i}.}

\partiel[15 min]{Questions}
\consigne{Répondez aux deux premières questions en français, aux deux dernières en anglais.}

\begin{enumerate}[label=\arabic*.,itemsep=2.5mm,leftmargin=*]
  \item Comment le poulain passait-il ses journées avant de pouvoir manger de l'herbe ?
    \lignes{2}
    \corr{Il se nourrissait du lait de sa mère, courait à côté d'elle le jour et se couchait
    tout près d'elle la nuit ; par temps chaud ils se tenaient à l'ombre près de la mare, par
    temps froid ils avaient un abri près du bosquet.}
  \item Quelle différence sa mère fait-elle entre lui et les autres poulains ?
    \lignes{2}
    \corr{Les autres sont des poulains de chevaux de trait : ils n'ont pas appris les bonnes
    manières. Lui est de bonne race — son père est connu, son grand-père a gagné la coupe à
    Newmarket, sa grand-mère avait le meilleur caractère.}
  \item \en{What exactly does his mother ask him to do, and not to do?}
    \lignes{3}
    \corr{To grow up gentle and good, never to learn bad ways, to do his work with a good
    will and to lift his feet well when he trots; and never to bite or kick, even in play.}
  \item \en{The horse tells his story long afterwards. Find two sentences that show it.}
    \lignes{2}
    \corr{\emph{I have never forgotten my mother's advice}, \emph{any horse I ever knew},
    \emph{Her name was Duchess, but he often called her Pet}. Le narrateur parle depuis un
    présent d'où il se souvient.}
\end{enumerate}

\note[\emph{Used to} et \emph{would}]{\emph{We used to stand by the pond}, \emph{they would
frequently bite} : deux façons de dire une habitude passée. \emph{Used to} + base verbale
signifie \og autrefois nous nous tenions \fg{} ; \emph{would} + base verbale ne s'emploie que
pour une action répétée, jamais pour un état. Les deux se traduisent par un imparfait.}

%% ===================================================================
\section{Version}
\consigne{Traduisez le passage en français. C'est la toute première page du roman.}

\begin{version}
The first place that I can well remember was a large pleasant
\gl[a meadow]{meadow}{un pré} with a pond of clear water in it. Some shady trees
\gl[to lean]{leaned}{se pencher} over it, and \gl[a rush]{rushes}{un jonc} and water-lilies
grew at the deep end. Over the hedge on one side we looked into a plowed field, and on the
other we looked over a gate at our master's house, which stood by the roadside; at the top of
the meadow was a grove of \gl[a fir tree]{fir trees}{un sapin}, and at the bottom a running
\gl[a brook]{brook}{un ruisseau} \gl[to overhang]{overhung}{surplomber} by a steep bank.
\end{version}
\source{\emph{Ibid.}}

\lignes{10}

\corr{\textbf{Proposition de traduction.} Le premier endroit dont je me souvienne bien était
un grand pré agréable, avec en son milieu une mare d'eau claire. Quelques arbres ombreux se
penchaient au-dessus, et des joncs et des nénuphars poussaient à l'extrémité profonde.
Par-dessus la haie, d'un côté, nous avions vue sur un champ labouré ; de l'autre, nous
regardions par-dessus une barrière la maison de notre maître, qui se tenait au bord de la
route ; en haut du pré il y avait un bosquet de sapins, et en bas un ruisseau vif que
surplombait une berge escarpée.}

\note[Choix de traduction 1]{\emph{The first place that I can well remember} : \emph{well}
porte sur \emph{remember}, \og dont je me souvienne bien \fg. Le subjonctif français est
d'usage après un superlatif ou un premier.}

\note[Choix de traduction 2]{\emph{at the top of the meadow was a grove} : l'anglais place le
complément de lieu en tête et inverse le sujet. Le français peut faire de même
(\og en haut du pré il y avait \fg) : c'est un des rares endroits où les deux langues
s'accordent.}

\note[Choix de traduction 3]{\emph{a running brook overhung by a steep bank} : \emph{running}
et \emph{overhung} sont deux participes, l'un actif, l'autre passif. Le français distingue
mieux en passant au relatif : \og un ruisseau vif que surplombait une berge \fg.}

%% ===================================================================
\partiel{Lexique à mémoriser}
\begin{multicols}{2}\small
\begin{itemize}[label=--,itemsep=0.8mm,leftmargin=*]
  \item \textbf{a meadow} — un pré
  \item \textbf{a pond} — une mare
  \item \textbf{shady} — ombragé
  \item \textbf{a hedge} — une haie \textit{(fiche 1)}
  \item \textbf{a gate} — une barrière, un portail
  \item \textbf{steep} — escarpé, raide
  \item \textbf{a colt} — un poulain
  \item \textbf{to kick} — ruer, donner un coup de pied
  \item \textbf{to bite} — mordre
  \item \textbf{manners} — les bonnes manières
  \item \textbf{advice} — les conseils \textit{(indénombrable)}
  \item \textbf{wise} — sage, avisé
\end{itemize}
\end{multicols}

\end{document}
